% Exercises for Chapter 3 (Arbitrage and static replication). Statements only: solutions
% are kept privately. Notation as in the notes: B is the price of the zero-coupon bond
% paying 1 at T, S_0 the stock price, C(K) and P(K) call and put prices with maturity T.

Unless an exercise says otherwise, the market is frictionless, the stock pays no dividend,
and the market is free of arbitrage in the sense of \noteref{def:arb-arbitrage}. Strikes,
prices and payoffs are in dollars.

\begin{exercise}[Building payoffs from calls]\label{ex:arb-payoffs}
\easy Use \noteref{lem:arb-piecewise}.
\begin{enumerate}[label=(\alph*)]
  \item Replicate the \emph{straddle}, which pays $|S_T-K|$ ($K>0$), with bonds, shares and
    calls.
    Replicate it again with one call and one put, and check with put--call parity that the
    two replications cost the same.
  \item A \emph{collar} pays $\min\big(\max(S_T,90),120\big)$. Replicate it with bonds and
    calls, and again with one share, one put and one call. Compare the costs.
  \item A payoff equals $10$ for $S_T\le80$, decreases linearly to $0$ at $S_T=90$, increases
    linearly to $5$ at $S_T=100$, and equals $5$ for $S_T\ge100$. Replicate it and write its
    price in terms of $B$ and call prices. Without computing anything, why is this price
    nonnegative?
  \item A \emph{lookback call} with strike $K$ pays $(M_T-K)^+$, where $M_T$ is the highest
    price of the stock between $0$ and $T$. Assuming it is traded and the market remains free
    of arbitrage, compare its price with $C(K)$.
\end{enumerate}
\end{exercise}

\begin{exercise}[A known dividend]\label{ex:arb-dividend}
\easy The stock now pays a known cash dividend $\delta>0$ at a date $t_1\in(0,T)$: whoever
holds the share at $t_1$ receives $\delta$, and a short seller must pay $\delta$ to the
lender of the share. Besides the bond maturing at $T$ (price $B$), a zero-coupon bond paying
$1$ at $t_1$ is traded, at price $B_1$. Portfolios are still static, but their payoffs now
consist of a cash flow at $t_1$ and one at $T$. You may use \noteref{prop:arb-dominance} for
such portfolios, with ``$V\ge\tilde V$'' meaning that the inequality holds at both dates.
\begin{enumerate}[label=(\alph*)]
  \item Show that the forward price for delivery at $T$ is $F=(S_0-\delta B_1)/B$.
  \item Show that $C(K)-P(K)=S_0-\delta B_1-BK$.
  \item Show that $(S_0-\delta B_1-BK)^+\le C(K)\le S_0-\delta B_1$.
  \item In \noteref{rem:arb-strict}, the lower bound for a call without dividends is at least
    the exercise value $S_0-K$ when $B\le1$. Show that this fails here: find conditions under
    which $S_0-\delta B_1-BK<S_0-K$ although $B_1\le1$ and $B\le1$.
\end{enumerate}
\end{exercise}

\begin{exercise}[A trainee's quote sheet]\label{ex:arb-quotes}
\medium The stock trades at $S_0=50$, the bond paying $1$ at $T=1$ at $B=0.96$, and a
trainee's program produces the following option prices for maturity $T$.
\begin{center}
\begin{tabular}{lccc}
  \toprule
  strike $K$ & $45$ & $50$ & $55$\\
  \midrule
  call & $7.90$ & $5.30$ & $2.40$\\
  put & $2.90$ & $4.60$ & $7.40$\\
  \bottomrule
\end{tabular}
\end{center}
For each arbitrage below, give the portfolio, the cash collected at time $0$, and the payoff
at time $T$ in every state.
\begin{enumerate}[label=(\alph*)]
  \item Check the calls alone against \noteref{prop:arb-bounds} and \noteref{prop:arb-strike},
    and exploit every violation you find.
  \item Check the puts alone in the same way.
  \item Check put--call parity at each strike and exploit every violation.
  \item A \emph{box spread} is long the $45/55$ call spread and short the $45/55$ put spread:
    long a call at $45$, short a call at $55$, short a put at $45$, long a put at $55$. Show
    that it pays $10$ in every state. What should it cost? What does it cost at the quoted
    prices, and which discount factor does that imply?
  \item The trainee made two mistakes: one call price is a typo, and the puts were computed
    from parity with a wrong discount factor. Using (a)--(d), identify both, and correct the
    call price.
\end{enumerate}
\end{exercise}

\begin{exercise}[What no-arbitrage leaves open]\label{ex:arb-range}
\easy Fix a strike $K>0$, and suppose that the market made of the bond, the stock and the
call with strike $K$ is free of arbitrage.
\begin{enumerate}[label=(\alph*)]
  \item Show that $\Prob(S_T>0)>0$, and deduce that $C(K)<S_0$. Show that $C(K)>0$ if
    $\Prob(S_T>K)>0$, and that $C(K)>S_0-BK$ if $\Prob(S_T<K)>0$.
  \item Suppose $S_0>BK$. Show that $C(K)=S_0-BK$ is possible only if $S_T\ge K$ almost
    surely, and that in that case the law of one price forces $C(K)=S_0-BK$.
  \item Let $c$ be a number with $(S_0-BK)^+<c<S_0$. Find $x>K$, a probability $\Prob$ under
    which $S_T$ takes the two values $0$ and $x$, each with positive probability, and a pricing
    measure (\noteref{def:arb-pricing-measure}) under which the call costs $c$. What does this say
    about the bounds of \noteref{prop:arb-bounds}?
  \item In your model of (c), what happens to $x$ and to $\Q(S_T=x)$ as $c$ increases to
    $S_0$? Interpret.
\end{enumerate}
\end{exercise}

\begin{exercise}[Digital options]\label{ex:arb-digital}
\medium A \emph{digital call} with strike $K>0$ pays $1$ if $S_T>K$ and $0$ otherwise; a
\emph{digital put} pays $1$ if $S_T\le K$. Their prices are $D_c(K)$ and $D_p(K)$. Digitals
and calls of every strike are traded, together with the bond and the stock.
\begin{enumerate}[label=(\alph*)]
  \item Show that $D_c(K)+D_p(K)=B$ and $0\le D_c(K)\le B$. With $B=0.98$, a dealer quotes a
    digital call at $1.01$: build an arbitrage.
  \item For $0<h<K$, show that
    \[
      \frac{C(K)-C(K+h)}{h}\le D_c(K)\le\frac{C(K-h)-C(K)}{h}.
    \]
  \item Deduce that $-\partial^+C(K)\le D_c(K)\le-\partial^-C(K)$. You may use that a convex
    function on an open interval has one-sided derivatives at every point, and that its
    difference quotients $\big(C(K+h)-C(K)\big)/h$ decrease to $\partial^+C(K)$ as
    $h\downarrow0$, and $\big(C(K)-C(K-h)\big)/h$ increase to $\partial^-C(K)$.
  \item Deduce that $D_c(K)=-C'(K)$ when $C$ is differentiable at $K$. If there is a pricing
    measure $\Q$, express $D_c(K)$ in terms of $\Q$ and compare with
    \noteref{prop:arb-breeden-litzenberger}.
  \item A trader has sold a digital call and wants to hold calls so that she never loses at
    $T$, whatever happens. Which portfolio of calls does (b) suggest, what does it cost, and
    what happens to this cost as $h\downarrow0$?
\end{enumerate}
\end{exercise}

\begin{exercise}[The log contract and the variance]\label{ex:arb-log}
\medium Let $\Q$ be a pricing measure for a market in which calls and puts of every strike
$K>0$ are traded, with $\Q(S_T>0)=1$, and let $F=S_0/B$. Assume that all the integrals of
option prices that appear below are finite, so that \noteref{cor:arb-strip} applies.
\begin{enumerate}[label=(\alph*)]
  \item Show that the claim paying $-\log(S_T/F)$ (the \emph{log contract}) has price
    \[
      L=\int_0^F\frac{P(K)}{K^2}\,\dd K+\int_F^\infty\frac{C(K)}{K^2}\,\dd K .
    \]
  \item Let $f:(0,\infty)\to\R$ be convex and twice continuously differentiable. Show, using
    \noteref{cor:arb-strip} and the sign of option prices only, that the claim paying
    $f(S_T)$ costs at least $Bf(F)$. Which classical inequality have you proved, and under
    which measure?
  \item Show that the claim paying $S_T^2$ costs
    $BF^2+2\int_0^FP(K)\,\dd K+2\int_F^\infty C(K)\,\dd K$, and deduce a formula for the
    variance of $S_T$ under $\Q$ in terms of option prices.
  \item Suppose that under $\Q$, $\log S_T$ is normal with variance $v^2$. Show that its mean
    must be $\log F-v^2/2$, and that $L=Bv^2/2$. (This is why the log contract is the
    building block of variance swaps, in \cref{ch:15-volatility}.)
\end{enumerate}
\end{exercise}

\begin{exercise}[Reading a density off rounded quotes]\label{ex:arb-density}
\medium This is a computational exercise. Take $S_0=100$, $B=0.97$, $F=S_0/B$, and suppose
that the market's call prices are
\[
  C(K)=S_0\,N(d_+)-BK\,N(d_-),\qquad d_\pm=\frac{\log(F/K)\pm v^2/2}{v},\qquad v=0.25,
\]
where $N$ is the standard normal distribution function. (This is the Black--Scholes formula of
\cref{ch:06-black-scholes-formula}; you may take it as given.) These prices come from a
pricing measure under which $S_T$ has the density
\[
  q_T(K)=\frac{1}{Kv\sqrt{2\pi}}\exp\Big(-\frac{\big(\log(K/F)+v^2/2\big)^2}{2v^2}\Big),
  \qquad K>0 .
\]
\begin{enumerate}[label=(\alph*)]
  \item At $K=100$ and $K=130$, compute the butterfly estimate
    \[
      \hat q_h(K)=\frac{C(K-h)-2C(K)+C(K+h)}{Bh^2}
    \]
    for $h\in\{20,10,5,2,1,0.5\}$ and compare it with $q_T(K)$. How does the error behave as
    $h\to0$? Explain the rate.
  \item Real quotes are rounded to the cent. Repeat (a) with the prices rounded to two
    decimals. What happens for small $h$, and why? Give a rough bound for the error caused by
    rounding, and choose a sensible $h$.
  \item Take $h=0.1$ at $K=130$: the rounded quotes give $\hat q_h(K)<0$. Which relation of
    the notes do they violate? Is there an arbitrage at the rounded prices, and would you
    expect to be able to trade it?
\end{enumerate}
\end{exercise}

\begin{exercise}[An arbitrage that depends on the model]\label{ex:arb-weak}
\hard Let $S_0=100$ and $B=1$, and suppose that calls with strikes $100$ and $110$ are both
quoted at $5$. In this exercise a \emph{model} is any probability $\Prob$ for $S_T\ge0$;
the prices are fixed.
\begin{enumerate}[label=(\alph*)]
  \item Check that the quotes satisfy every inequality of \noteref{prop:arb-bounds} and
    \noteref{prop:arb-strike}.
  \item Show that, whatever the model, the market made of the bond, the stock and the two
    calls has an arbitrage. (Distinguish the models with $\Prob(S_T>100)>0$ from the others.)
  \item Show that each of the two portfolios you used in (b) fails to be an arbitrage in some
    model.
  \item Show that no portfolio of the bond, the stock and the two calls is an arbitrage in
    every model. \emph{Hint:} if $\theta$ were one, first use models in which $S_T$ is
    constant to show that its payoff $V(x)$ is nonnegative for every $x\ge0$. Then, for
    $x>110$, use the model in which $S_T\in\{0,100,x\}$ with probabilities chosen so that
    the expected payoffs of the bond, the stock and the call with strike $100$ equal their
    prices, and let $x\to\infty$.
\end{enumerate}
\end{exercise}
